\section{Model Fixtures}

The current fixture suite contains three model audit styles.

\begin{center}
\small
\begin{tabularx}{\textwidth}{lXX}
\toprule
Fixture & Positive claim & Negative pair \\
\midrule
\texttt{colorado-rla} & A Colorado-style comparison audit record has a 20 digit public seed, SHA-256 sampler, margin metadata, stopping method, risk estimate, and status. & \texttt{bad-colorado-rla} rejects a seed that is not 20 decimal digits. \\
\texttt{california-rla} & A California-style public audit record names ballot manifest format, audit software id, public source URL, margin, and status. & \texttt{bad-california-rla} rejects non-public source metadata. \\
\texttt{manual-audit} & A plain manual audit hand count matches canvassed machine totals within zero-vote tolerance. & \texttt{bad-manual-audit} catches a hand-count drift hidden behind declared pass. \\
\bottomrule
\end{tabularx}
\end{center}

These fixtures are not real state certifications. They are controlled examples
that exercise different evidence styles. Colorado's public RLA materials
describe ballot manifests, public seed generation, and audit software sampling
\citep{colorado-rla-faq}; Rule 25 also names CVR and summary reconciliation
requirements for comparison audits \citep{colorado-rule-25}. California's RLA
regulations emphasize an accurate ballot manifest and public software
algorithm/source-code disclosure \citep{california-rla-regulations}. The manual
audit fixture covers the more ordinary case: compare hand-count totals to the
canvassed machine totals and escalate when they disagree beyond tolerance.
